\documentclass{zjureport}
\special{dvipdfmx:config z 0} % 取消PDF压缩，加快速度，最终版本生成的时候最好把这句话注释掉

% =============================================
% Part 0 Edit the info
% =============================================

\major{信息工程}
\name{黄鹤翔}
\partner{}
\title{本科实验报告}
\stuid{3230106231}
\college{信息与电子工程学院}
\date{\today}
\lab{东四227}
\course{场波实验}
\instructor{王子立}
\grades{}
\expname{CST仿真实验}
\exptype{仿真型}

\usepackage{listings}
\usepackage{xcolor}

\lstset{
 columns=fixed,       
 numbers=left,                                        % 在左侧显示行号
 numberstyle=\tiny\color{gray},                       % 设定行号格式
 frame=none,                                          % 不显示背景边框
 backgroundcolor=\color[RGB]{245,245,244},            % 设定背景颜色
 keywordstyle=\color[RGB]{40,40,255},                 % 设定关键字颜色
 numberstyle=\footnotesize\color{darkgray},           
 commentstyle=\it\color[RGB]{0,96,96},                % 设置代码注释的格式
 stringstyle=\rmfamily\slshape\color[RGB]{128,0,0},   % 设置字符串格式
 showstringspaces=false,                              % 不显示字符串中的空格
 breaklines=true,                                     % 对过长的代码自动换行
 language=verilog,                                     % 设置语言
}


\begin{document}
% =============================================
% Part 1 Header
% =============================================
\makecover

% \makecontent

\makeheader
% =============================================
% Part 2 Main document
% =============================================

% \begin{figure}[H] %H为当前位置，!htb为忽略美学标准，htbp为浮动图形
%     \centering %图片居中
%     \includegraphics[width=0.8\textwidth]{fig1-SR830structure.jpg}
%     \caption{锁相放大器SR830结构图} %最终文档中希望显示的图片标题
% \end{figure}
\section{实验目的}
\begin{itemize}
  \item 熟悉CST软件的使用方法
  \item 了解矩形波导馈电的角锥喇叭天线的结构、方向性等
  \item 掌握矩形波导馈电角锥喇叭天线的理论分析
  \item 通过矩形波导馈电角锥喇叭天线CST仿真观察方向图和喇叭中电场等情况
\end{itemize}

\section{实验内容}
角锥喇叭天线尺寸数据如下：$D_H = 80mm,D_E = 38mm,L = 80mm$，波导尺寸
$a = 22.86mm,b = 10.16mm,Lambda = 29.1mm$，铜壁厚$t = 1mm$，用CST软件
建模并仿真，观察方向图和喇叭中电场等情况。

\section{实验原理}
\begin{center}
  \includegraphics[width = 0.8\textwidth]{原理1.png}
  \includegraphics[width = 0.8\textwidth]{原理2.png}
  \includegraphics[width = 0.8\textwidth]{原理3.png}
  \includegraphics[width = 0.8\textwidth]{原理4.png}
\end{center}

\section{实验过程}
\subsection{建模}
\subsubsection{创建长方体}
创建的长方体参数与效果图如下图所示：\par
\begin{figure}[H]
  \begin{minipage}{0.45\textwidth}
    \centering
    \includegraphics[width=0.8\textwidth]{长方体创建2.png}
    \caption{长方体参数}
  \end{minipage}
  \begin{minipage}{0.45\textwidth}
    \centering
    \includegraphics[width=0.8\textwidth]{长方体创建1.png}
    \caption{长方体效果图}
  \end{minipage}
\end{figure}

\subsubsection{创建矩形}
创建的矩形参数与效果图如下图所示：\par
\begin{figure}[H]
  \begin{minipage}{0.45\textwidth}
    \centering
    \includegraphics[width=0.8\textwidth]{矩形创建1.png}
    \caption{矩形参数}
  \end{minipage}
  \begin{minipage}{0.45\textwidth}
    \centering
    \includegraphics[width=0.8\textwidth]{矩形创建2.png}
    \caption{矩形效果图}
  \end{minipage}
\end{figure}

\subsubsection{矩形平移}
矩形平移参数与效果图如下图所示：\par
\begin{figure}[H]
  \begin{minipage}{0.45\textwidth}
    \centering
    \includegraphics[width=0.8\textwidth]{矩形平移1.png}
    \caption{矩形平移参数}
  \end{minipage}
  \begin{minipage}{0.45\textwidth}
    \centering
    \includegraphics[width=0.8\textwidth]{矩形平移2.png}
    \caption{矩形平移效果图}
  \end{minipage}
\end{figure}

\subsubsection{侧壁创建与合并}
侧壁创建与合并的过程如下图所示：\par
\begin{figure}[H]
  \begin{minipage}{0.45\textwidth}
    \centering
    \includegraphics[width=0.8\textwidth]{侧壁创建.png}
    \caption{侧壁创建}
  \end{minipage}
  \begin{minipage}{0.45\textwidth}
    \centering
    \includegraphics[width=0.8\textwidth]{整体合并.png}
    \caption{整体合并}
  \end{minipage}
\end{figure}

\subsubsection{掏空}
掏空的过程与效果图如下图所示：\par
\begin{figure}[H]
  \begin{minipage}{0.45\textwidth}
    \centering
    \includegraphics[width=0.8\textwidth]{掏空1.png}
    \caption{掏空过程}
  \end{minipage}
  \begin{minipage}{0.45\textwidth}
    \centering
    \includegraphics[width=0.8\textwidth]{掏空2.png}
    \caption{掏空参数}
  \end{minipage}
\end{figure}
\begin{figure}[H]
  \centering
  \includegraphics[width=0.8\textwidth]{掏空3.png}
  \caption{掏空效果图}
\end{figure}

\subsection{仿真}
\subsubsection{仿真设置}
分别设置仿真频率、仿真背景与仿真边界条件
\begin{figure}[H]
  \begin{minipage}{0.32\textwidth}
    \centering
    \includegraphics[width=0.8\textwidth]{仿真设置1.png}
    \caption{仿真频率设置} 
  \end{minipage}
  \begin{minipage}{0.32\textwidth}
    \centering
    \includegraphics[width=0.8\textwidth]{仿真设置2.png}
    \caption{仿真背景设置}
  \end{minipage}
  \begin{minipage}{0.32\textwidth}
    \centering
    \includegraphics[width=0.8\textwidth]{仿真设置3.png}
    \caption{仿真边界条件设置}
  \end{minipage}
\end{figure}

\subsubsection{激励源与监视器设置}
\begin{figure}[H]
  \begin{minipage}{0.45\textwidth}
    \centering
    \includegraphics[width=0.8\textwidth]{激励源设置.png}
    \caption{激励源设置}
  \end{minipage}
  \begin{minipage}{0.45\textwidth}
    \centering
    \includegraphics[width=0.8\textwidth]{监视器设置.png}
    \caption{监视器设置}
  \end{minipage}
\end{figure}

\subsubsection{模式分析}
首先只计算端口模式不执行整个时域仿真，预先了解模式分布\par
\begin{figure}[H]
  \centering
  \includegraphics[width=0.8\textwidth]{仿真1.png}
  \caption{模式分布预了解}
\end{figure}
查看模式仿真分析报告，发现只有模式一的频率低于最高频率12.4GHz\par
\begin{figure}[H]
  \centering
  \includegraphics[width=0.8\textwidth]{仿真2.png}
  \caption{模式分析报告}
\end{figure}
选择端口1与模式1进行完整仿真
\begin{figure}[H]
  \centering
  \includegraphics[width=0.8\textwidth]{仿真3.png}
  \caption{端口1与模式1仿真}
\end{figure}

\section{实验结果与分析}
\subsection{喇叭天线S11}
\begin{figure}[H]
  \centering
  \includegraphics[width=0.8\textwidth]{s11.png}
  \caption{喇叭天线S11}
\end{figure}
由$S_{11}$曲线可以看出天线输入端口的反射系数在
工作频率8.2GHz到12.4GHz之间小于-20dB，说明天线
在工作频段内能较好地工作

\subsection{驻波比特性曲线}
\begin{figure}[H]
  \centering
  \includegraphics[width=0.8\textwidth]{驻波比.png}
  \caption{驻波比特性曲线}
\end{figure}
由驻波比特性曲线可以看出天线在频率8.2GHz到8.5GHz之间
反射系数较小，能够较好地工作，在其他频率段反射系数基本大于1.1，
工作状态不是很好

\subsection{方向图}
\begin{figure}[H]
  \begin{minipage}{0.45\textwidth}
    \centering
    \includegraphics[width=0.8\textwidth]{方向图1.png}
    \caption{3D方向图}
  \end{minipage}
  \begin{minipage}{0.45\textwidth}
    \centering
    \includegraphics[width=0.8\textwidth]{方向图2.png}
    \caption{2D方向图}
  \end{minipage}
\end{figure}

\subsection{增益图}
\begin{figure}[H]
  \begin{minipage}{0.45\textwidth}
    \centering
    \includegraphics[width=0.8\textwidth]{增益1.png}
    \caption{3D增益图}
  \end{minipage}
  \begin{minipage}{0.45\textwidth}
    \centering
    \includegraphics[width=0.8\textwidth]{增益2.png}
    \caption{2D增益图}
  \end{minipage}
\end{figure}

从仿真结果来看，该矩形波导馈电地角锥喇叭天线
地主瓣方向为$\varphi = 0,\theta = 0$，主瓣宽度为
$20.9^{\circ}$，主瓣的最大增益为$15.6dB$。\par
\subsubsection{理论增益计算}
\begin{align*}
  R_1 = \frac{L *D_H}{D_H-a} = 112.01mm\\
  R_2 = \frac{L *D_E}{D_E-b} = 200.79mm\\
  L_H = \sqrt{R_1^2 + \frac{D_H^2}{4}} = 118.94mm\\
  L_E = \sqrt{R_2^2 + \frac{D_E^2}{4}} = 201.69mm\\
  A = D_H*D_E = 3040mm^2\\
  S_H = \frac{A^2}{8\lambda L_H} = 333.76mm^2\\
  S_E = \frac{A^2}{8\lambda L_E} = 196.82mm^2\\
  \alpha = 8S_H = 2670.08mm^2\\
  \beta = 8S_E = 1574.56mm^2\\
\end{align*}
读表可得$\Delta G_H = 0.3292,\Delta G_E = 0.6960$\\
可以计算得到天线增益为：
\begin{align*}
  G &= 10.8+10\lg(\frac{D_H*D_E}{\lambda^2}) - \Delta G_H - \Delta G_E\\
    &= 10.8 + 10\lg(\frac{80*38}{29.1^2}) - 0.3292 - 0.6960\\
    &= 15.326 dB\\
\end{align*}
理论值与仿真值大致相同

\subsection{电场情况}
\begin{figure}[H]
  \begin{minipage}{0.45\textwidth}
    \centering
    \includegraphics[width=0.8\textwidth]{电场1.png}
    \caption{电场分布图1}
  \end{minipage}
  \begin{minipage}{0.45\textwidth}
    \centering
    \includegraphics[width=0.8\textwidth]{电场2.png}
    \caption{电场分布图2}
  \end{minipage}
\end{figure}
从图中可以看出，两边的电场关于yz平面对称，
且离平面越远电场越弱；波导电场最强，在喇叭中传播
越远电场越弱；喇叭天线内电场强，外部电场弱。

\section{实验结论}
从仿真结果来看，该矩形波导馈电地角锥喇叭天线
地主瓣方向为$\varphi = 0,\theta = 0$，主瓣宽度为
$20.9^{\circ}$，主瓣的最大增益为$15.6dB$，与理论结果接近。
同时天线输入端口的反射系数在
工作频率8.2GHz到12.4GHz之间小于-20dB，说明天线
在工作频段内能较好地工作。

\section{收获与体会}
本次实验通过CST软件对矩形波导馈电的角锥喇叭天线进行建模与仿真，
让我学会了CST的基本操作，熟悉了仿真的基本流程，相信对我之后的课程都会很有帮助。
\end{document}